% !PRINTED 523
\weberpage{523}

adjungiren, worin
\weq{5}{D = + a^2b^2 + 18\,abc - 4a^3c - 4b^3 - 27c^2}
die Discriminante der Gleichung (3) ist, so reducirt sich die Gruppe
auf (2).

In dem K\"orper $\Omega'$, der durch diese Adjunction aus $\Omega$
entsteht, kann jede Wurzel rational durch jede andere ausgedr\"uckt
werden, denn es ist
\[
\begin{gathered}
\alpha_1 + \alpha_2 = a - \alpha\\
\alpha_1 - \alpha_2 = \dfrac{\sqrt{D}}{f'(\alpha)},
\end{gathered}
\]
also
\weq{6}{\begin{gathered}
2\alpha_1 = a - \alpha + \dfrac{\sqrt{D}}{f'(\alpha)}\\
2\alpha_2 = a - \alpha - \dfrac{\sqrt{D}}{f'(\alpha)},
\end{gathered}}
und hierin k\"onnen die Vertauschungen $\pi$, $\pi^2$ ausgef\"uhrt
werden. Die cubische Gleichung ist also nach Adjunction von $\sqrt{D}$
ihre eigene \gesperrt{Galois'sche} Resolvente.

Alles dies bleibt g\"ultig, wenn der Rationalit\"atsbereich irgend
ein specieller K\"orper $\Omega$ ist, in dem $a$, $b$, $c$ und
$\sqrt{D}$ enthalten sind, wenn nicht $f(x)$ selbst in $\Omega$
reducibel ist, also eine rationale Wurzel hat. Denn ausser sich
selbst und der cyklischen Gruppe (2) hat die Gruppe (1) nur noch
intransitive Theiler, n\"amlich die Einheitsgruppe und drei Gruppen
vom Typus (1, 2).

Wollen wir die cubische Gleichung nun aufl\"osen, d.\,h. auf eine
reine cubische Gleichung zur\"uckf\"uhren, so m\"ussen wir eine
Function $v$ der Wurzeln suchen, die zwar nicht selbst, deren Cubus
aber die cyklische Permutation $\pi$ gestattet. Eine solche Function
kann aber nur existiren, wenn die dritten Einheitswurzeln, oder, was
dasselbe ist, $\sqrt{-3}$ dem K\"orper $\Omega$ adjungirt wird; denn
$v$ muss durch Anwendung von $\pi$ eine dritte Einheitswurzel als
Factor erhalten.

Durch diese Adjunction kann eine Reduction der Gruppe nicht
eintreten, weil die Gruppe (2) ausser der Einheitsgruppe keinen
Theiler hat und die Reduction auf die Einheitsgruppe nicht durch
eine quadratische, sondern nur durch eine cubische Gleichung
geschehen kann (\S.\,157). Wir setzen also
\[
\varepsilon = \dfrac{-1+\sqrt{-3}}{2}, \qquad
\varepsilon^2 = \dfrac{-1-\sqrt{-3}}{3}
\]
% [sic] Der Druck gibt im Nenner von \varepsilon^2 die Zahl 3; algebraisch wäre 2 erforderlich. Der Druckwert bleibt quellentreu erhalten.
%CONT>>
